\documentclass[11pt]{article}
\usepackage[margin=1in]{geometry}
\usepackage{amsmath,amssymb,amsthm}
\usepackage[colorlinks=true,linkcolor=blue,urlcolor=blue]{hyperref}
\newtheorem{theorem}{Theorem}
\newtheorem{proposition}[theorem]{Proposition}
\newtheorem{remark}[theorem]{Remark}
\newcommand{\one}{\boldsymbol 1}
\newcommand{\tr}{\operatorname{tr}}
\newcommand{\R}{\mathbb R}
\title{A bounded-degree weighted obstruction to gap-free balanced-partition stability}
\author{Independent theorem audit; research working note}
\date{November 15, 2026}
\begin{document}
\maketitle

\paragraph{Scope and status.}
Every vertex has ordinary unit mass; all Laplacians below are
combinatorial, and all vector norms and projectors use the usual
Euclidean inner product. The hosts in Theorem~\ref{thm:counter} are
simple, connected, bipartite, and have combinatorial maximum degree
four. They have \emph{weighted} edges, including large internal
conductances. No unweighted-host assertion is made. The edited edge
has conductance exactly one before the update, and that one edge is
doubled. The theorem is an explicit counterexample, not a priority
claim. Proposition~\ref{prop:ridge} records a classical resolvent and
determinant consequence, not a new matrix inequality.

\paragraph{Big-if target that fails.}
Could there be a dimension-independent function $\omega(t)\to0$ as
$t\downarrow0$ such that doubling one existing edge of leverage $\tau$
changes a unique optimal unit-mass bisection on at most
$\omega(\tau)n$ vertices, or changes its unique Fiedler projector in
operator norm by at most $\omega(\tau)$? The answer is no, even in the
bounded-degree graph class just specified, and even when the entire
fractional relative correction tends to zero in operator norm.
The obstruction is to labels and hard spectral projectors; objective
values can remain very close.

\begin{theorem}[One unit edge; vanishing leverage; half the labels change]
\label{thm:counter}
For every integer $r\ge2$, put $n=4r$. There is a weighted graph $G_r$
on $n$ unit-mass vertices, with underlying graph $C_4\mathbin\square P_r$,
and an existing unit-conductance edge $e_r$, with these properties.
Let $\widehat L$ be its combinatorial Laplacian and
$\widehat L'=\widehat L+bb^T$, where $b$ is the incidence vector of $e_r$.
\begin{enumerate}
\item The initial and updated Fiedler eigenvalues are simple, and
corresponding unit eigenvectors $x,y$ have every coordinate equal to
$1/\sqrt n$ or $-1/\sqrt n$. They are orthogonal.
\item Their positive-coordinate sets $S_x,S_y$ are the unique minimum
bisections of $G_r,G_r'$, respectively, up to complementation. Both
have cardinality $n/2$, and
\[
 \min\{|S_x\triangle S_y|,|S_x\triangle S_y^c|\}=n/2.
\]
\item The edit leverage satisfies
\[
 0<\tau_r=b^T\widehat L^\dagger b<\frac4n.
\]
The old and new optimal bisection weights are $n/2$ and $n/2+1/2$.
The old cut has updated weight $n/2+1$.
\item For every fixed $0<\alpha<1$, the relative fractional correction
\[
 C_r=(\widehat L^\dagger)^{\alpha/2}
       (\widehat L'{}^\alpha-\widehat L^\alpha)
       (\widehat L^\dagger)^{\alpha/2}
\]
obeys
\[
 \|C_r\|\le(1+\tau_r)^\alpha-1
 \le\alpha\tau_r\longrightarrow0,
\]
whereas the Fiedler projectors satisfy
$\|xx^T-yy^T\|=1$ and $\|xx^T-yy^T\|_F^2=2$.
\end{enumerate}
\end{theorem}

\begin{proof}
Set $\delta=1/n$ and $M=nr^2$. Label vertices by $(a,j)$,
$a\in\{1,2,3,4\}$ and $1\le j\le r$. At each $j$, put square edges
\[
 (1,j)(2,j),\ (3,j)(4,j)\quad\text{of conductance }1+\delta,
\]
\[
 (1,j)(3,j),\ (2,j)(4,j)\quad\text{of conductance }1.
\]
For each fixed $a$, join $(a,j)$ and $(a,j+1)$ by an edge of
conductance $2M$. This is the Cartesian product of a square and a
path as an unweighted underlying graph. It is simple, connected,
and has maximum degree at most four. It is bipartite: add the parity
of $j$ to a square bipartition, for which $\{1,4\}$ and $\{2,3\}$
are the two classes.

We analyze $A=\widehat L/2$; multiplication by two changes neither
eigenvectors nor minimizing cuts, and the edit becomes
$B=A+\frac12bb^T$. Let $K_\delta$ be the four-vertex weighted square
with conductances $(1+\delta)/2$ on $12,34$ and $1/2$ on $13,24$.
Then
\[
 A=L(K_\delta)\otimes I_r+M I_4\otimes L(P_r).
\]
The three unit Hadamard vectors
\[
 q_x=\tfrac12(1,1,-1,-1)^T,\quad
 q_y=\tfrac12(1,-1,1,-1)^T,\quad
 q_z=\tfrac12(1,-1,-1,1)^T
\]
are eigenvectors of $L(K_\delta)$ with eigenvalues
$1,1+\delta,2+\delta$, respectively. Define
\[
 x=q_x\otimes\one_r/\sqrt r,\qquad
 y=q_y\otimes\one_r/\sqrt r,\qquad
 z=q_z\otimes\one_r/\sqrt r.
\]
All remaining positive eigenvalues of $A$ are at least
$M\nu_1$, where
\[
 \nu_1=2-2\cos(\pi/r)=4\sin^2(\pi/(2r))\ge4/r^2.
\]
Here the last inequality is $\sin u\ge2u/\pi$ on $[0,\pi/2]$.
Consequently $M\nu_1\ge4n$. The Fiedler vector of $A$, and hence
of $\widehat L$, is uniquely $x$ up to sign.

Fix any $j_0$ and double $e_r=(1,j_0)(3,j_0)$. Thus
$b=(e_1-e_3)\otimes e_{j_0}$. We have
\[
 b^Ty=0,\qquad b^Tx=b^Tz=\beta:=2/\sqrt n,
 \qquad b=\beta x+\beta z+h,
 \quad \|h\|^2=2-8/n,
\]
where $h$ lies in the spectral subspace of $A$ above $4n$.
Thus $y$ remains an eigenvector of $B$ at $t=1+\delta$.
We show it is uniquely the lowest positive eigenvector.

On $\{\one,y\}^\perp$, the original spectrum has its unique lowest
eigenvalue at $1$, its next at $c=2+\delta$, and all others at least
$4n$. For $s\in(1,c)$ the rank-one secular function is
\[
 f(s)=1+\frac12\left(
    \frac{\beta^2}{1-s}+\frac{\beta^2}{c-s}
    +\langle h,(A-sI)^{-1}h\rangle\right).
\]
It is strictly increasing from $-\infty$ to $+\infty$; its unique
root $\mu$ is the smallest eigenvalue of the updated restriction.
All other eigenvalues of that restriction are at least $c$, by
rank-one interlacing. At $t=1+\delta$,
\[
 f(t)
 =-1+\frac2n+\frac12\langle h,(A-tI)^{-1}h\rangle
 \le-1+\frac2n+\frac1{4n-1-1/n}<0\qquad(n\ge8).
\]
Therefore $\mu>1+\delta$. It follows that $y$ is the unique Fiedler
vector of $B$, and of $\widehat L'$, up to sign. The Fiedler
eigenvalues before and after the original scaling are
$2$ and $2(1+\delta)$.

For any bisection $S$, its normalized sign vector
$f_S=(2\one_S-\one)/\sqrt n$ has norm one and is orthogonal to
constants. Moreover
\[
 f_S^T\widehat L f_S=\frac4n w_G(\partial S).
\]
The Rayleigh variational principle bounds this energy below by $2$,
with equality only for $f_S=\pm x$. Thus $S_x$ is the unique minimum
bisection up to complementation. The same argument using the simple
updated eigenvalue $2(1+\delta)$ shows uniqueness of $S_y$ after
the edit. Equivalently their weights follow directly from the square
rails: $w_G(\partial S_x)=n/2$ and
$w_G(\partial S_y)=(n/2)(1+\delta)=n/2+1/2$; the doubled edge
crosses $S_x$ and lies inside $S_y$.
The Hadamard signs give the claimed symmetric difference.

For leverage, the scaling gives
\[
 \tau_r=\frac12 b^TA^\dagger b
 =\frac2n\left(1+\frac1{2+\delta}\right)
        +\frac12\langle h,A^\dagger h\rangle
 \le\frac2n\left(1+\frac1{2+\delta}\right)+\frac1{4n}
 <\frac{13}{4n}<\frac4n.
\]
For an exact scalar formula, let $(\nu_k,\phi_k)_{k=0}^{r-1}$ be
an orthonormal spectral resolution of $L(P_r)$, with
$\nu_0=0$ and $\phi_0=\one_r/\sqrt r$. Then
\[
 \tau_r=\frac12\sum_{k=0}^{r-1}|\phi_k(j_0)|^2
 \left(\frac1{1+M\nu_k}+
       \frac1{2+\delta+M\nu_k}\right).
\]
All inverses on constants are interpreted as zero; $b$ is orthogonal
to constants, so this convention causes no ambiguity.

Finally the positive fractional response has exact logarithmic budget
$\tr\log(I+C_r)=\alpha\log(1+\tau_r)$. Its largest eigenvalue is
therefore at most $(1+\tau_r)^\alpha-1$, which is at most
$\alpha\tau_r$ by scalar concavity. Orthogonality of $x,y$ gives
the stated projector distances. Positive fractional powers preserve
the eigenvectors and their order, so the same Fiedler obstruction
holds for $\widehat L^\alpha$ and $\widehat L'{}^\alpha$.
\end{proof}

\begin{remark}[Why the weighted qualifier matters]
The large path conductances enforce rigid equal-size groups, and
$1+1/n$ supplies the half-unit difference between competing cut
weights. In a fully unweighted graph, a single unit-edge addition
or doubling cannot replace one unique minimum bisection by a
different unique minimum bisection: every competing cut originally
has integer weight at least one larger, whereas the old optimum
increases by at most one. A tie may appear, or two edits may cause
a strict change. This elementary integrality observation does not
protect unique Fiedler eigenvectors in unweighted graphs.
\end{remark}

\paragraph{A scale-averaged replacement, with classical status.}
Hard projectors and optimizer labels fail in Theorem~\ref{thm:counter}.
Soft ridge filters admit a clean gap-free budget. This follows
directly from Sherman--Morrison and logarithmic differentiation of
a perturbation determinant; it should not be sold as a breakthrough.

\begin{proposition}[Exact logarithmic ridge-response budget]
\label{prop:ridge}
Let $A=L(G)$ for a finite connected weighted graph with unit masses,
$B=A+dbb^T$ with $d>0$, and $\tau=db^TA^\dagger b$. For $s>0$ put
\[
 F_A(s)=s(A+sI)^{-1},\quad F_B(s)=s(B+sI)^{-1},\quad
 D_s=F_A(s)-F_B(s).
\]
Then $D_s\succeq0$, $\operatorname{rank}D_s=1$, $D_s\preceq I$, and
\[
 \int_0^\infty\tr D_s\,\frac{ds}{s}
 =\int_0^\infty\|D_s\|\,\frac{ds}{s}
 =\log(1+\tau).
\]
In particular,
\[
 \int_0^\infty\|D_s\|_F^2\,\frac{ds}{s}
 \le\frac{\tau}{1+\tau}\log(1+\tau),
\qquad
 \int_{\{s:\|D_s\|>\eta\}}\frac{ds}{s}
 \le\frac{\log(1+\tau)}\eta\quad(\eta>0).
\]
\end{proposition}
\begin{proof}
Put $v_s=(A+sI)^{-1}b$ and $r(s)=b^T(A+sI)^{-1}b$.
Sherman--Morrison gives
\[
 D_s=\frac{sd}{1+dr(s)}v_sv_s^T.
\]
This is positive rank one, and $0\preceq D_s\preceq F_A(s)\preceq I$.
Furthermore
\[
 \tr D_s=\frac{sd\|v_s\|^2}{1+dr(s)}
        =-s\frac{d}{ds}\log(1+dr(s)).
\]
As $s\downarrow0$, $dr(s)\to\tau$, since $b\perp\one$;
as $s\to\infty$, $dr(s)\to0$. Integrating proves the identity.
The trace of a positive rank-one matrix is its operator norm. Also
$s(A+sI)^{-2}\preceq(A+sI)^{-1}$, so
$\tr D_s\le dr(s)/(1+dr(s))\le\tau/(1+\tau)$.
Consequently $\|D_s\|_F^2=(\tr D_s)^2
\le[\tau/(1+\tau)]\tr D_s$, proving the stronger small-leverage bound.
The final bound is Markov's inequality for logarithmic scale measure.
\end{proof}

\paragraph{What remains open after these audits.}
A viable positive theorem must either change the output to a smooth
filter, allow an explicitly ambiguous family of cuts, or assume a
partition-specific objective or rounding margin. A claim for hard
Fiedler projectors based only on leverage or the fractional log-budget
is false. A claim for exact balanced-cut labels based only on those
quantities is false for weighted hosts even at maximum degree four.
The relative fractional factor theorem provides useful uniform energy
certificates, but extracting these statements from it does not improve
the ordinary spectral partition approximation ratio, since
$L^\alpha$ has the same eigenvectors as $L$.
\end{document}
